% TOP_PROBLEM: {"id": 30006773, "title": "Inhomogeneous Landau–Coulomb Global Regularity (HST All-Moment Data)", "category_id": 9, "category": {"id": 9, "name": "pde", "display_name": "Partial Differential Equations", "description": "PDEs and their applications in physics and geometry.", "slug": "pde", "order_index": 9, "created_at": "2026-07-31T15:26:25.671Z"}, "status": "open"}
% FIRST_POSTED: 2026-09-27
% !TeX program = lualatex
% Compile twice: lualatex 30006773.tex
% All references and prose are embedded; no BibTeX database or external inputs.
\documentclass[11pt,a4paper]{article}
\usepackage[margin=24mm,headheight=14pt]{geometry}
\usepackage{fontspec}
\setmainfont{Latin Modern Roman}
\setsansfont{Latin Modern Sans}
\setmonofont{Latin Modern Mono}
\usepackage{amsmath,amssymb,mathtools}
\usepackage{unicode-math}
\setmathfont{Latin Modern Math}
\usepackage{microtype}
\usepackage{enumitem,xurl,needspace}
\usepackage[dvipsnames]{xcolor}
\usepackage{hyperref}
\hypersetup{unicode=true,colorlinks=true,linkcolor=MidnightBlue,urlcolor=MidnightBlue,
 pdftitle={500 Mathematical Problem Statements},pdfauthor={Source-annotated compilation},bookmarksnumbered=true}
\usepackage{bookmark}
\usepackage{tocloft}
\setlength{\cftsecnumwidth}{3em}
\cftsetpnumwidth{2.5em}
\cftsetrmarg{4em}
\usepackage{titlesec}
\titleformat{\section}{\Large\bfseries\raggedright}{\thesection}{0.7em}{}
\usepackage{fancyhdr}
\pagestyle{fancy}\fancyhf{}
\fancyhead[L]{\small\sffamily Mathematical problem statements}
\fancyhead[R]{\small Source edition 22}
\fancyfoot[C]{\thepage}
\setlength{\parindent}{0pt}
\setlength{\parskip}{5pt plus 1pt}
\setlength{\emergencystretch}{3em}
\setcounter{tocdepth}{1}
\setcounter{secnumdepth}{1}
\setlist[itemize]{leftmargin=3em,itemsep=5pt,topsep=4pt,parsep=0pt}
\allowdisplaybreaks
\newcommand{\N}{\mathbb{N}}
\newcommand{\Z}{\mathbb{Z}}
\newcommand{\Q}{\mathbb{Q}}
\newcommand{\R}{\mathbb{R}}
\newcommand{\C}{\mathbb{C}}
\newcommand{\F}{\mathbb{F}}
\newcommand{\A}{\mathbb{A}}
\newcommand{\PP}{\mathbb P}
\newcommand{\E}{\mathbb{E}}
\newcommand{\eps}{\varepsilon}
\newcommand{\dd}{\,\mathrm d}
\newcommand{\sref}[2]{\hyperlink{source:#1:#2}{[S#2]}}
\newcommand{\eref}[2]{\hyperlink{extra:#1:#2}{[E#2]}}
% Turn the next switch off for a shorter reading copy. The quotations preserve
% every exactTarget field, including qualifications omitted from a short summary.
\newif\ifshowcatalogue
\showcataloguetrue
\newcommand{\cataloguescope}[1]{\ifshowcatalogue
 \paragraph{Exact scope in the supplied catalogue.}
 \begin{quote}\small #1\end{quote}\fi}

% Standalone notebook wrapper; original problem section retained verbatim.
\numberwithin{equation}{section}
\hypersetup{pdftitle={Research attempt -- problem 30006773},
 pdfauthor={Research notebook; source compilation retained}}
\fancyhead[L]{\small\sffamily Research notebook}
\fancyhead[R]{\small\sffamily Problem 30006773 | 27 September 2026}
\begin{document}
\setcounter{section}{0}
% ================================================================
% problemId: 30006773
% Canonical problem ID: 30006773
% exactTarget SHA-256: 934833b6da68511281cb0364caac141eb1297f0c48a66c035743783bc608f298
\clearpage
\section*{\texorpdfstring{Inhomogeneous Landau–Coulomb Global Regularity (HST All-Moment Data)}{Inhomogeneous Landau–Coulomb Global Regularity (HST All-Moment Data)}}\label{30006773}
\subsection{Definitions and mathematical statement}
Write $\langle v\rangle=(1+|v|^2)^{1/2}$. The inhomogeneous Landau equation transports a nonnegative particle density $f(t,x,v)$ by
$(\partial_t+v\cdot\nabla_x)f=Q(f,f)$ with the Coulomb matrix
$A(z)=a|z|^{-1}(I-z\otimes z/|z|^2)$, $a>0$. The collision operator and the all-polynomial-moment, global H\"older, and uniformly positive local-ball hypotheses are retained exactly in the catalogue quotation below. Let $T_E$ be the maximal continuation time of the classical branch selected in the HST source. The problem is $T_E=\infty$ for every admitted initial datum. The source continuation quantity is
\[
 \Psi(T)=\mathop{\mathrm{ess\,sup}}_{t\le T,x}\int_{\R^3}f(t,x,v)\dd v
 +\mathop{\mathrm{ess\,sup}}_{t\le T,x,v}f(t,x,v).
\]
Establishing its finiteness on every finite time interval is a possible continuation route, not an assumption granting the conclusion. Total mass over all space need not be finite.
\par\smallskip\textit{Formulation and concept references:} \sref{30006773}{1}.
\cataloguescope{Fix a positive physical constant a. On [0,infinity) times R\textasciicircum{}3\_x times R\textasciicircum{}3\_v consider the classical Landau equation (partial\_t+v dot grad\_x)f=Q(f,f), where Q(f,f)=div\_v integral\_\{R\textasciicircum{}3\} A(v-w)[f(t,x,w) grad\_v f(t,x,v)-f(t,x,v) grad\_w f(t,x,w)]dw and A(z)=a |z|\textasciicircum{}(-1)(I-z tensor z/|z|\textasciicircum{}2) for z!=0. For every nonnegative initial datum f\_0 on R\textasciicircum{}6 with f\_0 in the global Holder space C\textasciicircum{}alpha for some 0<alpha<1, sup\_\{x,v\}<v>\textasciicircum{}k f\_0(x,v)<infinity for every k>=0, and constants delta,r\_0,R>0 such that for every x there exists |v\_x|<R with f\_0>=delta on the Euclidean R\textasciicircum{}6 ball B\_\{r\_0\}(x,v\_x), does the classical solution branch constructed in HST Theorems 1.2–1.4 have maximal continuation time T\_E=infinity? Here <v>=(1+|v|\textasciicircum{}2)\textasciicircum{}(1/2); the branch matches f\_0 as t decreases to zero, is smooth for positive time, and has all polynomial velocity weights bounded on every finite time slab of its existence. The exact question is global extension of that branch, with the sufficiently decaying uniformly continuous kinetic weak uniqueness convention in HST Section 5, not uniqueness among all generalized weak solutions. The HST Coulomb continuation quantity is Psi(T)=esssup\_\{0<=t<=T,x\} integral f(t,x,v)dv + esssup\_\{0<=t<=T,x,v\} f(t,x,v). Finite Psi on every finite time slab is the associated continuation criterion, not an extra hypothesis presumed true. No smallness near equilibrium, finite total spatial mass or uniform bound across different initial data is imposed.}
\subsection{Short English statement}
Does the specified classical Landau–Coulomb solution continue for all time for every initial density in the stated all-moment class, without assuming smallness or bounded continuation quantities in advance?
% BEGIN RESEARCH ADDITION ATTEMPT-179
\subsection{Research attempt}

\paragraph{Current frontier.}
The HST paper is Henderson--Snelson--Tarfulea's local-solution theorem;
Theorems 1.2--1.4 and Section 5 give the continuation and uniqueness
conventions used in the original entry \sref{30006773}{1}. For Coulomb
interactions the continuation criterion contains both the uniform local
velocity mass and $\|f\|_\infty$. Guillen--Silvestre prove global regularity
and monotonicity of Fisher information for the \emph{spatially homogeneous}
Landau equation \eref{30006773}{1}, Theorems 1.1--1.2. That important result
cannot be read as a theorem for the inhomogeneous branch here, and the
initial hypotheses of that branch are not replaced by finite total mass
in all of phase space.

\paragraph{Attack route.}
A maximum-principle route would need the diffusion to offset the local
quadratic reaction. A moment route would need to control local spatial
concentration, not only conserved total quantities. A transfer of the
homogeneous Fisher-information argument must handle commutators with
transport. We examine all three at the level of explicit identities and
construct admissible smooth test data that break the simplest maximum
estimate. The goal is to determine which proposed a priori input could
actually close the HST continuation criterion.

\paragraph{Attempt: the Coulomb reaction has a fixed positive sign.}
Let $\mathcal A[f](v)=\int A(v-w)f(w)\,dw$, with $x,t$ suppressed.
Since $A(z)=aD^2|z|$ in three dimensions,
\[
 b_i(z):=\partial_jA_{ij}(z)=-2a\frac{z_i}{|z|^3},
 \qquad \partial_i b_i=-8\pi a\delta_0
\]
in distributions. The coefficient of $\delta_0$ follows from the flux of
$z/|z|^3$ through a sphere, which is $4\pi$.
Integrating the $w$ derivative in the collision operator and expanding
the $v$ divergence cancels its two first-derivative terms. Therefore, for
smooth rapidly decaying velocity profiles,
\begin{equation}\label{eq179-nondiv}
 (\partial_t+v\cdot\nabla_x)f
       =\mathcal A[f]:D_v^2f+8\pi a f^2.
\end{equation}
This normalization uses the physical constant in the question. A convention
absorbing $8\pi a$ into time would not justify dropping it in this display.

At an attained maximum in $(x,v)$, diffusion is nonpositive and transport
vanishes, so the elementary maximum comparison gives only the Riccati
upper inequality $M'\le8\pi a M^2$. Where the usual localized supremum
comparison is justified this yields
\[
       M(t)\le\frac{M(0)}{1-8\pi a M(0)t}
          \quad\text{for }t<(8\pi aM(0))^{-1}.
\]
This is a short-time upper bound, not proof of blow-up and not a global
bound. On the noncompact spatial domain a localization argument is needed
to promote pointwise maximum reasoning to a supremum; the countertest below
uses spatially homogeneous data and an actual velocity plateau, avoiding
that separate issue entirely.

\paragraph{Attempt: admissible plateaus defeat a naive damping lemma.}
Choose $G(v)>0$ smooth, rapidly decreasing, and constant equal to $c>0$
on $|v|\le1$. Choose $\chi\in C_c^\infty(\R^3)$ with
$0\le\chi\le1$ and $\chi=1$ near the origin. For small $\eps>0$ let
\begin{equation}\label{eq179-spike}
 f_{0,\eps}(x,v)=G(v)+\eps^{-1}\chi(v/\eps),
\end{equation}
independent of $x$. Each datum satisfies the stated global H\"older and
all-polynomial-weight bounds, with constants allowed to depend on $\eps$.
The positive local-ball hypothesis holds uniformly in $x$ by taking
$v_x=0$ and a small fixed radius; the baseline supplies a fixed lower bound.
Thus these are legitimate data for the chosen class, not excluded spikes.

At $(x,0)$ the velocity Hessian and the spatial gradient vanish. For these
smooth data the equation at initial time, or its smooth right-hand limit,
gives exactly
\begin{equation}\label{eq179-positive}
       \partial_t f_{\eps}(0,x,0)
                         =8\pi a(c+\eps^{-1})^2>0.
\end{equation}
The spike adds only $O(\eps^2)$ to velocity mass, $O(\eps^4)$ to energy,
and $O(\eps^2\log(1/\eps))$ to the positive entropy contribution.
Every fixed integrated polynomial moment remains bounded in this family.
The entropy estimates follow by separating the support of the spike and
using its volume $O(\eps^3)$ and height $O(\eps^{-1})$; the baseline has
finite entropy. Therefore bounded mass, energy, and entropy alone cannot
force a universal pointwise inequality with an instantaneous damping term
that dominates the quadratic growth at all high flat maxima.

This does not contradict \eref{30006773}{1}: diffusion may act immediately off
the plateau and subsequently control the evolution by a nonlocal mechanism.
The construction is not evidence of a solution blowing up. It refutes only
a proposed instantaneous maximum-damping estimate derived from those
coarse quantities. Uniform H\"older or pointwise weighted bounds across
this family have deliberately not been asserted.

\paragraph{Attempt: the exact local-mass interpolation.}
Let $M=\|f\|_{L^\infty_{x,v}}$ and
$F_k=\|\langle v\rangle^k f\|_{L^\infty_{x,v}}$ for a fixed $k>3$.
For $R\ge1$, split the velocity integral into a ball and its complement:
\[
 \sup_x\int f(x,v)\,dv\le CMR^3+C_kF_kR^{3-k}.
\]
If $M>0$, then $F_k\ge M$ and $R=(F_k/M)^{1/k}\ge1$ is admissible.
It follows that
\begin{equation}\label{eq179-interpolation}
 \sup_x\int f(x,v)\,dv
       \le C_k M^{1-3/k}F_k^{3/k}.
\end{equation}
For $M=0$ the statement is immediate. Thus bounds on $M$ and one such
weighted supremum up to a putative finite maximal time would bound $\Psi$
and invoke HST continuation. But boundedness of every weight on each
\emph{compact subinterval} of existence does not provide a bound uniform
as that maximal time is approached. Assuming the latter would insert the
missing estimate into the hypotheses.

Velocity integration of the equation yields
$\partial_t\rho+\nabla_x\cdot j=0$, not a maximum principle for $\rho$.
The flux $j=\int vf\,dv$ has no sign in this equation. Hence even a
conserved global spatial mass, when finite in a smaller data class, would
not bound the spatial supremum in \eqref{eq179-interpolation}. For the
actual class it need not be finite at all.

\paragraph{Attempt: transport obstructs the homogeneous monotonicity.}
On a smooth positive rapidly decaying phase-space test solution where the
integrals are finite, set
$I_v(f)=\int f|\nabla_v\log f|^2\,dx\,dv$.
The transport part alone contributes
\begin{equation}\label{eq179-cross}
 \left.\frac{d}{dt}I_v(f)\right|_{\partial_t f=-v\cdot\nabla_xf}
       =-2\int f\,\nabla_v\log f\cdot\nabla_x\log f\,dx\,dv.
\end{equation}
Indeed $D_t\log f=0$ for free transport, but
$D_t\partial_{v_i}\log f=-\partial_{x_i}\log f$ because
$[\partial_{v_i},D_t]=\partial_{x_i}$; integration of the transport
divergence has no boundary term under these decay assumptions.
The right-hand side has no fixed sign. Thus homogeneous monotonicity of
$I_v$ does not transfer simply by adding transport. A hypocoercive
combination involving spatial derivatives is a plausible repair, but it
would require new estimates that control those derivatives and the local
mass without assuming the desired continuation bound.

\paragraph{Stress test.}
Equation \eqref{eq179-cross} is used only as a finite-integral diagnostic,
not as a globally defined functional for all HST data. The admissible
spatially homogeneous plateau family has infinite total spatial mass and
is nonetheless allowed by the target. Mixing these two tests into one
unjustified global integration argument would be an error.

The distributional Coulomb divergence, the positive sign of its reaction,
and the distinction between integrated moments and weighted supremum
bounds have all been tracked. Neither the Riccati comparison nor the
plateau growth proves finite-time blow-up. Conversely, the homogeneous
global theorem does not supply the missing estimates for the spatial flux
or the commutator term.

\paragraph{Outcome.}
\textbf{Outcome: Exploratory approach with unresolved gap.}

The attempt derives the exact reaction identity, an admissible countertest
to a naive maximum-damping argument, a local-mass interpolation bound, and
the transport obstruction to direct Fisher-information monotonicity.
These computations identify specific missing estimates rather than a new
global regularity theorem. What remains is a uniform finite-time bound
closing $\Psi$ for every admitted inhomogeneous datum. No blow-up solution
or complete continuation argument is claimed.
% END RESEARCH ADDITION ATTEMPT-179
\subsection{Sources}
\begin{itemize}\raggedright
\item[\textnormal{[S1]}]\hypertarget{source:30006773:1}{} Inhomogeneous Landau–Coulomb Global Regularity (HST All-Moment Data); exact editorial selection using the primary formulations and hypotheses recorded in the source receipts.
\url{https://www.numdam.org/item/10.1016/j.anihpc.2020.04.004.pdf}.
{\small\textit{Catalogue formulation source.}}
% BEGIN RESEARCH ADDITION SOURCES-179
\item[\textnormal{[E1]}]\hypertarget{extra:30006773:1}{} Nestor Guillen and
Luis Silvestre, \emph{The Landau equation does not blow up},
Acta Mathematica 234(2) (2025); arXiv:2311.09420v2,
Theorems 1.1--1.2 and Section 11.
\url{https://arxiv.org/html/2311.09420v2}.
{\small\textit{Consulted 27 September 2026. The cited theorem is
spatially homogeneous; the inhomogeneous transport term is not covered.}}
% END RESEARCH ADDITION SOURCES-179
\end{itemize}

\end{document}
